\section{Discussion}
\label{sec:discussion}

\subsection{The engine-ridge framework}
\label{sec:engine-ridge}
Once the compute engine is a free variable, each precision has its own ridge
(Figure~\ref{fig:ridge}). On B300 the ridge is $\ai^\star = \text{peak} /
8\tbs$: $\approx$312 for FP16 (2.5\,PF), 625 for FP8 (5\,PF), and 1875 for NVFP4
(15\,PF). This reframes the precision levers. FP8 MLA-latent storage gives
$\ai \approx 470 < 625$, so even with the tensor-core engine the kernel is
HBM-leaning, consistent with v12's measured 46.3\% HBM $>$ 35.7\% compute at
scale. NVFP4 pushes \ai\ higher but its 15\,PF peak lifts the ridge to 1875, so
the FP4 compute peak is \emph{doubly} unreachable for decode: the workload is
HBM-bound before it is compute-bound. The engine-ridge picture is why a bigger
FLOP peak does not, by itself, help a memory-lean decode workload.

\begin{figure}[t]
  \centering
  \includegraphics[width=\columnwidth]{v12-engine-ridge-comparison.pdf}
  \caption{Engine-ridge comparison on B300. Each precision's ridge scales with
  its tensor-core peak (FP16 $\approx$312, FP8 625, NVFP4 1875 FLOP/byte). Decode
  arithmetic intensities land left of every ridge, so NVFP4's 15\,PF peak
  overshoots into the HBM-bound region and is unreachable for decode.}
  \label{fig:ridge}
\end{figure}

\subsection{Implications for serving systems}
\label{sec:serving}
The measurements separate two regimes cleanly. \emph{Single-stream,
moderate-context decode}, the common interactive case, leaves even the
production kernel at 1--2\% of peak; the lever there is not precision or algorithm
but \emph{work} (batching and context length that fill the tensor-core pipeline).
\emph{Long-context, large-batch serving} does reach the roofs (46\% HBM at
$B{=}64$, $K{=}524$K). Consequently the value of FP8/NVFP4 KV is
\emph{capacity and accuracy} (halving or quartering the cache, letting more
sequences or longer contexts fit), not decode latency: the FP8 latency win is
L2-resident-only and reverses under an L2 flush, and NVFP4 is latency-negative
past L2 (12--30\% slower than FP16). Practitioners should treat low-precision KV
as a capacity knob and reach for engine plus batch to convert decode work into
throughput.

\subsection{Does native FP4 compute help MLA decode? (Boundary test)}
\label{sec:stagea}
Because MLA decode meets the block-scaled NVFP4 tensor-core gate ($M{\ge}128$) by
construction, it is tempting to expect that native FP4 \emph{compute}
(15\,PF, $3\times$ FP8) would accelerate it. We pre-registered the prediction
``no'' and tested it directly on B300 (\texttt{sm\_103a}, CUDA 12.9, CUTLASS
v4.5.2): a sweep of the $M{=}128$ MLA-decode QK-GEMM over
$K \in \{256,512,576\}$, $N \in \{1\text{K}\ldots524\text{K}\}$, comparing FP4
(CUTLASS example 72a, block-scaled NVFP4$\to$bf16) against FP8 (cuBLAS E4M3 via
\texttt{torch.\_scaled\_mm}, matched bf16 output). Figure~\ref{fig:arm2} and the
data of record confirm the \textbf{kill}: both precisions are HBM-bound at every
$N$ ($\ai$ 100--176 $\ll$ the ridges 1875 / 625), reaching at most \textbf{5.8\%}
of the 15\,PF FP4 peak and \textbf{21.7\%} of the 5\,PF FP8 peak. Native FP4
\emph{compute} is therefore never the operative limit, and the $M{=}128$ packing
advantage is a red herring \emph{for decode}. The raw FP4/FP8 ratio flips with
$N$: FP4 leads $1.4$--$2.7\times$ at small $N$ (both tiny, latency-bound), while
FP8 wins at large $N$, the realistic long-context regime (e.g. $K{=}576$,
$N{=}131072$: FP8 1086 vs FP4 739\tflops). Neither direction is a compute
win; the gap is a bandwidth / operand-size effect, i.e. the same
capacity/storage lever as Sections~\ref{sec:results-decode}
and~\ref{sec:serving}, not tensor-core throughput.

\begin{figure}[t]
  \centering
  \includegraphics[width=\columnwidth]{v12-arm2-research-plan.pdf}
  \caption{The native-FP4-compute boundary test (Section~\ref{sec:stagea}). At the
  $M{=}128$ MLA-decode QK-GEMM shape, both FP4 and FP8 GEMMs are HBM-bound at
  every context length; FP4 tops out at 5.8\% of its 15\,PF peak. The result is a
  pre-registered negative: FP4/NVFP4 stays a capacity+accuracy lever, not a
  decode-compute lever.}
  \label{fig:arm2}
\end{figure}

This is a cross-harness comparison (a CUTLASS \emph{example} FP4 kernel versus a
tuned cuBLAS FP8 kernel), so the head-to-head ratio is not clean; but the kill
rests on the \emph{kernel-quality-independent} peak-fraction and roofline
($\ai \ll$ ridge caps FP4 far below its compute peak regardless of tuning), not
on the ratio. The consequence is that accuracy follow-ups for a fused FP4-compute
kernel are \emph{not} triggered, and FP4/NVFP4 remains a capacity$+$accuracy
lever. (We use ``Stage A'' for this GEMM-compute test; it is distinct from the
v10 asymmetric-precision \emph{accuracy} ablation, which we do not abbreviate the
same way.)

\subsection{Threats to validity}
\label{sec:threats}
We disclose the following.
\emph{Counter-free proxy:} the \%HBM proxy is \texttt{ncu}-validated only on a
root T4 (13.8\% vs 12.85\%); on B300 \texttt{ncu} was blocked
(unprivileged container) and the proxy carries as a throughput estimate, not a
counter measurement.
\emph{Clock variability:} Colab throttles; the load-bearing verdicts use
locked-clock root-T4 runs and idle-pinned Blackwell runs, but cross-kernel
$\mu$s/tok on Blackwell is not fully clock-controlled.
\emph{Single-seed accuracy:} FP8/NVFP4 RMSE figures are single-seed on
gpt2-small and on synthetic KV; they are our substrate numbers, not published
constants, and the $\sim$3.5$\times$ FP4-vs-FP8 latent gap should be read as such.
\emph{v5 was never benchmarked:} its $8\times$ floor drop is a prediction only;
we never present it as measured.
\emph{nsys provenance:} the v10 B300 schedule figures come from an off-notebook
\texttt{nsys} 2025.3.2 run (the committed cell ran 2025.1.3, which records an
empty \texttt{sm\_103} trace); the v11 schedule was re-captured with 2025.3.2 and
is non-empty in-repo.
\emph{No SOTA comparators run:} FlashMLA-class kernels are orders of magnitude
faster by construction; we make no wall-clock claim against them. The
contribution is the open methodology and characterization.
